\section{Related Work}
\label{sec:related}

\para{The dead salmon tradition in science.}
\citet{bennett2009neural} demonstrated that standard fMRI analysis applied to a dead Atlantic salmon produced statistically significant brain activations when multiple comparisons corrections were not applied.
This work catalyzed widespread adoption of family-wise error rate (FWER) and false discovery rate (FDR) corrections in neuroimaging.
\citet{eklund2016cluster} extended this finding at scale, showing that widely used fMRI software packages (SPM, FSL, AFNI) produced inflated false-positive rates up to 70\%, potentially affecting tens of thousands of published studies.
More recently, \citet{meloux2025dead} applied the dead salmon concept to AI interpretability, demonstrating that feature attribution methods, probing classifiers, and sparse autoencoders produce plausible but spurious explanations for randomly initialized neural networks.
Our work extends this tradition to \aiscientist systems: rather than testing a specific analytical tool, we test the end-to-end analytical reasoning of \llms acting as scientists.
Unlike prior dead salmon studies that focus on false positives from null data, we also provide a ``live salmon'' with known signal, enabling measurement of both false positive and false negative rates.

\para{Autonomous AI scientist systems.}
\citet{lu2024ai} introduced The AI Scientist, the first framework for fully automated scientific discovery, where \llms generate ideas, write code, run experiments, produce papers, and conduct peer review---all for approximately \$15 per paper.
Subsequent evaluations have identified significant limitations.
\citet{beel2025evaluating} found that 42\% of AI Scientist experiments failed due to coding errors, literature reviews were inadequate (median 5 citations), and numerical results were sometimes hallucinated.
\citet{luo2025more} identified four systematic failure modes: inappropriate benchmark selection, data leakage, metric misuse, and post-hoc selection bias.
\citet{zhu2025ai} argued that the fundamental bottleneck is implementation capability, not idea generation---\llms produce novel ideas but struggle to execute rigorous experiments.
Our work differs from these evaluations in a key respect: we provide a controlled ground truth.
Rather than auditing AI-generated papers for quality or plausibility, we measure whether \aiscientists recover a known data generating process, enabling precise quantification of what they find and what they miss.

\para{Benchmarking AI research agents.}
\citet{starace2025paperbench} introduced PaperBench, decomposing 20 ICML 2024 papers into 8,316 gradable sub-tasks and finding that the best agent achieved only 21\% replication score.
\citet{gu2024blade} proposed BLADE, a benchmark for evaluating agents on data-driven science with ground truth from expert data scientists.
\citet{si2024can} conducted a large-scale study with 100+ NLP researchers, finding that LLM-generated research ideas are rated as more novel but less feasible than human expert ideas.
Our live salmon benchmark complements these efforts by providing a simpler but more controlled evaluation: rather than measuring broad research capability, we isolate the analytical reasoning component and test it against a known answer.
This design enables precise diagnosis of \textit{what types of patterns} \aiscientists can and cannot detect.
